\scanpage{9}
Soit $C=\operatorname{Rep}G$ la $\otimes$-catégorie définie par $G$, et $S$ un espace analytique \uline{connexe}, $s\in S$. On veut classifier les $\otimes$-foncteurs
\[
F:C\longrightarrow\text{str. de Hodge \inserted{(mod. isogénie)} sur }S
\]
telles que\tikz[baseline=-.2ex,x=1ex,y=1ex]{\draw[line width=.55pt] (0,2.0)--(0,.13) .. controls (-.30,-.07) and (.32,-.16) .. (.17,.10) .. controls (.13,.21) and (.05,.19) .. (0,.13);}\inserted{\tikz[baseline=-.2ex,x=1ex,y=1ex]{\draw[line width=.42pt] (.10,2)--(0,.16) .. controls (-.26,-.02) and (.37,-.20) .. (.28,.11) .. controls (.23,.25) and (.08,.24) .. (0,.16);}\ a) degré total (rel. à $i$) $=\sum$ degrés partiels b) $\mathbb{Q}(-1)$ donne str. de Hodge de bidegré $(1,1)$, et} c) Pour \deleted{tout $s\in S$ [en \uline{un} $s\in S$, c'est kif kif]} \deleted{et} $V\in\operatorname{Ob}C$, \deleted{\unsure{polarisables en chaque point}} \inserted{$V$ de degré $n$, $\exists\,\varphi:V\otimes V\unsure{\rightharpoonup|\to}\mathbb{Q}(-n)$ tel que $F(\varphi)$ soit une polarisation [en un pt, en tt point, de $S$, c'est kif-kif].}
\margin{De tels objets s'appellent \uline{$G$-subordination sur $S$} (pas néc. \unsure{rigidifiée|rigidifiées}!)}
Composant avec le foncteur fibre en $s$, on trouve un foncteur fibre
\[
C\longrightarrow\operatorname{Modf}(\mathbb{Q}),
\]
et quitte à tordre $G$, on peut supposer qu'on se donné un isomorphisme \inserted{$\varphi$} entre ce dernier et le foncteur oubli.

\deleted{Alors} Soit $\pi=\pi_1(S,\unsure{\xi|\zeta|s})$. Alors $\pi$ opère \inserted{à g.}, sur le foncteur \deleted{\unsure{``fibre}} ``$\mathbb{Q}$-fibre en $s$'' sur la catégorie des str. de Hodge sur $S$, donc opère sur le foncteur composé avec $F$, d'où un premier invariant de $(F,\varphi)$ :
\[
\alpha:\pi=\pi_1(S,\unsure{\xi|\zeta})\longrightarrow G(\mathbb{Q})
\]
Soit d'autre part $\widetilde{S}$ le \unsure{revêtt|revêtmt} universel de $S$ \deleted{\unsure{pointé}} \unsure{marqué} sur $s$ par $\widetilde{s}$. Alors la donnée de $\varphi$ permet de considérer le foncteur $V\longmapsto F(V)\longmapsto F(W)\times_{\unsure{s|S}}\overline{S}$ comme une $G$-subordination de $\widetilde{S}$. Donc on trouve un
% archive sheet number: 8

